% eksperymentalne tłumaczenie na włoski

%

\section{Twierdzenie sinusów} % lang-pl
\index{twierdzenie!sinusów}% % lang-pl
\section{Teorema dei seni} % lang-it
\index{teorema!dei seni}% % lang-it

\begin{proposition}
[twierdzenie sinusów] % lang-pl
[teorema dei seni] % lang-it
    Niech trójkąt o bokach długości $a, b, c$ oraz kątach $\alpha$, $\beta$, $\gamma$ naprzeciwko tychże boków (w tej samej kolejności) będzie wpisany w okrąg o promieniu $R$. % lang-pl
    Wtedy % lang-pl
    Sia un triangolo con lati di lunghezze $a, b, c$ e angoli $\alpha$, $\beta$, $\gamma$ opposti a questi lati (nello stesso ordine), inscritto in un cerchio di raggio $R$. % lang-it
    Allora % lang-it
    \label{twierdzenie_sinusow}%
    \begin{equation}
        \frac{a}{\sin \alpha} = \frac{b}{\sin \beta} = \frac{c}{\sin \gamma} = 2R.
    \end{equation}
\end{proposition}

Ptolemeusz w II wieku będzie wiedzieć o równoważnej formie, że boki trójkąta są proporcjonalne do cięciw podwojonych kątów naprzeciwko. % lang-pl
Brahmagupta w VII wieku wyrazi promień okręgu opisanego jako iloczyn dwóch boków trójkąta podzielony przez podwojoną wysokość, stąd można już wyprowadzić twierdzenie. % lang-pl
Sferyczną wersję pierwszy poda któryś z~muzułmańskich matematyków X wieku: Abu-Mahmud al-Khujandi, Abu Nasr Mansur albo Abu al-Wafa. % lang-pl
Tolomeo nel II secolo conoscerà una forma equivalente, secondo cui i lati del triangolo sono proporzionali alle corde dei doppi angoli opposti. % lang-it
Brahmagupta nel VII secolo esprimerà il raggio del cerchio circoscritto come il prodotto di due lati del triangolo diviso per il doppio dell'altezza; da qui si può già dedurre il teorema. % lang-it
Una versione sferica sarà data per primo da uno dei matematici musulmani del X secolo: Abu-Mahmud al-Khujandi, Abu Nasr Mansur oppure Abu al-Wafa. % lang-it

Coxeter \cite[s. 28, 29]{coxeter_1967} wyprowadza twierdzenie sinusów z symetrii okręgu opisanego na trójkącie. % lang-pl
Hartshorne \cite[s. 185]{hartshorne_2000} podaje je jako ćwiczenie. % lang-pl
Bogdańska, Neugebauer \cite[s. 89]{neugebauer_2018} oprą się o fakt \ref{maybe_3_28}, że kąty wpisane oparte o ten sam łuk mają równe miary. % lang-pl
Coxeter \cite[pp. 28, 29]{coxeter_1967} deduce il teorema dei seni dalla simmetria del cerchio circoscritto al triangolo. % lang-it
Hartshorne \cite[p. 185]{hartshorne_2000} lo propone come esercizio. % lang-it
Bogdańska e Neugebauer \cite[p. 89]{neugebauer_2018} si baseranno sul fatto \ref{maybe_3_28}, che gli angoli inscritti insistenti sullo stesso arco hanno misure uguali. % lang-it

%